% ============================================================
%  Conclusion
% ============================================================

\section{Conclusion}%
\label{sec:conclusion}

The preceding sections have presented the algorithms, data structures, and
engineering decisions underlying \FractalShark{}, a GPU-accelerated
Mandelbrot set renderer capable of interactive exploration at zoom
depths exceeding $10^{10000}$.  \Cref{fig:system-architecture}
provides a high-level view of the pipeline.

\begin{figure}[!htbp]
\centering
\begin{tikzpicture}[
  every node/.style={font=\small},
  block/.style={draw, rounded corners, minimum height=0.65cm,
                minimum width=2.1cm, align=center},
  comp/.style={block, fill=blue!10},
  io/.style={block, fill=orange!10},
  infra/.style={block, fill=gray!12, dashed},
  arr/.style={-{Stealth[length=5pt]}, thick}
]

% --- Row 1: Input ---
\node[io] (params) {View Parameters};
\node[comp, right=0.5cm of params] (pzbb) {PointZoom-\\BBConverter};
\node[comp, right=0.5cm of pzbb] (deltas) {Per-pixel\\deltas};

\draw[arr] (params) -- (pzbb);
\draw[arr] (pzbb) -- (deltas);

% --- Row 2: Core computation ---
\node[comp, below=1.0cm of params] (reforbit) {Reference Orbit\\(RefOrbitCalc)};
\node[comp, right=0.5cm of reforbit] (la) {LA Setup\\(LAReference)};
\node[comp, right=0.5cm of la] (perturb) {Perturbation\\Kernels};

\draw[arr] (deltas.south) -- ++(0,-0.45) -| (perturb.north);
\draw[arr] (reforbit) -- (la);
\draw[arr] (la) -- (perturb);

% --- Row 3: Output ---
\node[comp, below=1.0cm of reforbit] (aa) {Palette\\Coloring};
\node[comp, right=0.5cm of aa] (color) {Supersample\\Averaging};
\node[io, right=0.5cm of color] (gl) {GL Display};

\draw[arr] (perturb.south) -- ++(0,-0.45) -| (aa.north);
\draw[arr] (aa) -- (color);
\draw[arr] (color) -- (gl);

% --- Side: GPU arithmetic ---
\node[infra, left=0.6cm of reforbit] (gpu) {GPU Arith.\\(HpSharkFloat)};
\draw[arr] (gpu) -- (reforbit);

% --- Side: Feature Finder ---
\node[infra, below=0.35cm of gpu] (ff) {Feature\\Finder};
\draw[arr, dashed] (ff.east) -- ++(0.3,0) |- ([yshift=-2pt]reforbit.south west);

% --- Side: Memory ---
\node[infra, below=0.35cm of ff] (mem) {Growable-\\Vector};
\draw[arr, dashed] (mem.east) -- ++(0.6,0) |- ([yshift=2pt]reforbit.south);

\end{tikzpicture}
\caption[End-to-end rendering pipeline]{High-level architecture of the
  \FractalShark{} rendering pipeline.  Solid arrows denote the primary
  data flow; dashed arrows indicate supporting subsystems.}
\label{fig:system-architecture}
\end{figure}

The exposition began with the direct rendering kernels that map pixels
to the complex parameter plane and evaluate the escape-time iteration
(\cref{sec:mapping}), followed by the numeric type hierarchy---from
native \code{float} through multi-limb \code{HpSharkFloat}---that
provides the precision each stage requires (\cref{sec:hdr-float}).
Building on this foundation, perturbation theory
(\cref{sec:perturbation-concept}) reduces per-pixel work from
arbitrary-precision arithmetic to lightweight floating-point deltas,
and approximation methods (\cref{sec:bilinear-approx,sec:la-v2-perturb})
skip large blocks of iterations using precomputed coefficients.
BLA~v1 uses purely linear maps; the current LA~v2 path retains a
nonlinear input operation before evaluating its staged maps.

The second half of the document addressed the computation of the
reference orbit that perturbation demands
(\cref{sec:ref-orbit-calc}), the NTT-based GPU arithmetic engine that
accelerates it (\cref{subsec:ref-orbit-gpu}), and the Feature Finder
that applies the same high-precision machinery to automatic detection
of periodic points (\cref{sec:feature-finder}).  Supporting
infrastructure---post-iteration processing and color mapping
(\cref{sec:misc-details}), together with the disk-backed growable
vectors that manage large reference orbits
(\cref{sec:disk_backed_growable_vectors})---completes the rendering
pipeline.

The main techniques address different costs. Perturbation reduces the precision
required for a pixel update; BLA and LA reduce how many updates must be executed.
Parallel reference arithmetic reduces the cost of constructing their shared input.
The resulting speed depends on the view, the validity of accepted skips, and the
balance between reference construction and pixel work.

Several design decisions merit emphasis. \FractalShark{} uses \emph{rebasing}
to change a perturbation's reference index while retaining its reconstructed state.
Rebasing reduces some cancellation problems, but does not certify the accuracy of
every subsequent operation or eliminate every possible rendering artifact.  Linear
approximation version~2 (LA~v2) was chosen over series approximation
for its adaptive staged skip structure, which can enable longer
effective jumps at deep zooms. Compared with BLA~v1, nonlinear
preparation and useful return boundaries can reduce the individual
updates required per pixel (\cref{sec:la-v2-motivation}), independently
of optional AT initialization. NTT-based multiplication was preferred over
Karatsuba for the GPU arithmetic backend, as the former maps naturally
to the highly parallel CUDA execution model.  Throughout the
codebase, C++ template instantiation over a fixed parameter-set
lattice allows the compiler to specialize kernels for each combination
of precision, sub-type, and algorithmic mode while keeping the source
generic.

Taken together, these techniques form a layered architecture in which
each subsystem can be understood, tested, and tuned independently.
\Cref{sec:future-directions} surveyed several alternative and
complementary techniques---including series approximation,
smooth coloring, distance estimation, and histogram-based coloring---that provide
context for possible extensions of the current design.

The final appendix (\cref{sec:gpu-png-encoding}) illustrates a supporting application
of GPU parallelism: reversible filtering and compression of the completed image.
Its role is output efficiency; the Mandelbrot and approximation mathematics do not
depend on the encoder.

\FractalShark{} remains under active development.  As new algorithmic
ideas emerge from the deep-zoom rendering community and as GPU
hardware continues to evolve, the architecture described here provides
a foundation for continued experimentation and improvement.
